% =====================================================================
% 第3章 ブランチ・PR編 — チームで安全に開発する
% 元原稿: 03_ブランチ・PR編.md
% =====================================================================
\chapter{ブランチ・PR編 —チームで安全に開発する—}

本章は第2章の続きです。第2章では、ローカルとGitHubをつないで \textbf{1本の流れ（\cmd{main}）} をアップロード／ダウンロードしました。本章では、\textbf{ブランチ}で作業を枝分かれさせ、\textbf{プルリクエスト（PR）}でレビューを挟み、\textbf{コンフリクト（競合）}を解決する——つまり「\textbf{複数人・複数タスクを安全に並行させる}」方法を扱います。

\begin{hintbox}
\textbf{この章の読み方}
\begin{itemize}[leftmargin=1.4em, nosep]
  \item 1人で使う分には第2章までで十分です。チームで開発する、または作業を安全に分けたくなったらここを読みます
  \item コマンドをド忘れした人 → 第5章「逆引きチートシート」へ
  \item エラーが出て困った人 → 第4章「トラブルシューティング編」へ
\end{itemize}
\end{hintbox}

% ---------------------------------------------------------------------
\section{なぜブランチが必要か}

これまでは \cmd{main} という1本の流れだけで作業してきました。しかし、\cmd{main} に直接編集を重ねていくと、こんな問題が起きます。

\begin{itemize}[leftmargin=1.4em]
  \item \textbf{完成していない変更が本流に混ざる}：試しに書いたコードや書きかけの原稿が、\cmd{main} を壊してしまう
  \item \textbf{複数人が同じ流れを触って衝突する}：AさんとBさんが同時に \cmd{main} を編集すると、\cmd{push} がぶつかる
  \item \textbf{お試しがやりにくい}：「うまくいかなかったら元に戻す」を安全に試せない
\end{itemize}

\textbf{ブランチ}を使うと、\cmd{main} から\textbf{枝分かれした作業用の流れ}を作れます。作業はその枝の上で行い、完成してから \cmd{main} に合流（マージ）させます。こうすれば、\textbf{\cmd{main} は常に「動く・完成した状態」に保たれます}（図\ref{fig:branch-merge}）。

\begin{numfig}{ブランチで枝分かれし、完成してから main に合流させる}{fig:branch-merge}
\begin{tikzpicture}
  \node[flownode, fill=gitGreen] (M0) {main};
  \node[flownode, fill=gitGreen, right=2.4cm of M0] (M1) {main};
  \node[flownode, fill=gitGreen, right=4.8cm of M1] (M2) {main\\（マージ後）};
  \node[flownode, fill=gitOrange, above=1.4cm of M1] (F1) {feature\\作業中};
  \node[flownode, fill=gitOrange, right=2.4cm of F1] (F2) {feature\\完成};
  \draw[flowarrow] (M0) -- (M1);
  \draw[flowarrow] (M1) -- (M2);
  \draw[flowarrow] (F1) -- (F2);
  \draw[flowarrow, dashed] (M0) to[out=70, in=200]
        node[edgelabel, midway, left] {枝分かれ} (F1);
  \draw[flowarrow, dashed] (F2) to[out=-20, in=110]
        node[edgelabel, midway, right] {マージ} (M2);
\end{tikzpicture}
\end{numfig}

% ---------------------------------------------------------------------
\section{ブランチとは}

\textbf{ブランチ（branch＝枝）}は、その名の通り\textbf{履歴の枝分かれ}です。ある時点の \cmd{main} からコピーを取り、そこから独立して作業を進める「もう1本の作業ライン」だと考えてください。

\begin{itemize}[leftmargin=1.4em]
  \item \cmd{main}（メインブランチ）… 完成した状態を保つ本流。通常ここは\textbf{直接編集しない}
  \item \cmd{feature/...} などの\textbf{作業ブランチ} … 1つのタスク専用の枝。ここで自由に編集・コミットする
\end{itemize}

各ブランチは独立しているので、作業ブランチで何をしても \cmd{main} には影響しません。そして「今どのブランチにいるか」を \textbf{\cmd{HEAD}} が指しています。\cmd{git switch} でブランチを切り替えると、作業ツリーの中身もそのブランチの内容に入れ替わります。

\begin{hintbox}
\textbf{1タスク＝1ブランチ}が基本です。「ログイン機能の追加」「第3章の加筆」のように、\textbf{1つの目的ごとにブランチを分ける}と、後述のPRやレビューがきれいにまとまります。
\end{hintbox}

\subsection{たとえ：倉庫で考えるブランチ}

第1章で、リポジトリを\textbf{書類を保管する「倉庫」}にたとえました。ブランチは、この倉庫の中に置かれた\textbf{棚}だと考えると分かりやすくなります。

\begin{center}
\begin{tabularx}{\linewidth}{@{}l X@{}}
\toprule
Git & 倉庫のたとえ \\
\midrule
\cmd{main} & 全員が共有する\textbf{正式な棚（本棚）}。常に「完成品」だけを並べる \\
作業ブランチ & 本棚の中身を\textbf{丸ごと写し取った、自分専用の作業棚} \\
そのブランチでの \cmd{commit} & 自分の作業棚に、書き直した書類の写しを並べていく \\
PR → マージ & 完成した書類を\textbf{本棚に納める申請と反映} \\
\bottomrule
\end{tabularx}
\end{center}

大事なのは、作業棚は\textbf{本棚の写しから始まる独立した棚}なので、そこで何をしても本棚（\cmd{main}）は汚れない、という点です。だから安心して書きかけの書類を置いたり、失敗した写しを捨てたりできます。

\subsection{複数人のときはどうなるか}

倉庫のたとえは、\textbf{複数人での作業}でこそ効いてきます。メンバーはそれぞれ、本棚から\textbf{自分専用の作業棚を写し取って}並行して作業します。互いの作業棚は独立しているので、\textbf{他の人の手元を邪魔しません}（図\ref{fig:multi-branch}）。

\begin{numfig}{各メンバーが本棚から作業棚を写し取り、完成したら本棚へ納める}{fig:multi-branch}
\begin{tikzpicture}
  \node[flownode, fill=gitGreen, text width=2.6cm] (MAIN)
        {本棚（main）\\完成品の共有棚};
  \node[flownode, fill=gitOrange, text width=3cm, above right=1.9cm and 5cm of MAIN] (A)
        {Aさんの作業棚\\feature/login};
  \node[flownode, fill=gitPurple, text width=3cm, below right=1.9cm and 5cm of MAIN] (B)
        {Bさんの作業棚\\feature/report};
  % Aさん：写し取る（破線・上の弧）／PR→マージ（実線・下の弧）
  \draw[flowarrow, dashed] (MAIN.north east) to[out=55, in=210]
        node[edgelabel, pos=0.62, above left=-1pt] {写し取る} (A.west);
  \draw[flowarrow] (A.south west) to[out=250, in=30]
        node[edgelabel, pos=0.4, below right=-1pt] {PR→マージ} (MAIN.east);
  % Bさん：写し取る（破線・下の弧）／PR→マージ（実線・上の弧）
  \draw[flowarrow, dashed] (MAIN.south east) to[out=-55, in=150]
        node[edgelabel, pos=0.62, below left=-1pt] {写し取る} (B.west);
  \draw[flowarrow] (B.north west) to[out=110, in=-30]
        node[edgelabel, pos=0.4, above right=-1pt] {PR→マージ} (MAIN.east);
\end{tikzpicture}
\end{numfig}

完成したら、各自が自分の作業棚の書類を順番に本棚へ納めて（マージして）いきます。ここで、AさんとBさんが「同じ書類の同じ箇所」を別々に書き換えていた場合、\textbf{後から本棚へ納める人のときに「どちらの内容を正とするか決められない」状態になります。} これが後述の\textbf{コンフリクト（競合）}です。

\begin{hintbox}
逆に言えば、\textbf{触る書類（ファイル）や箇所をタスクごとに分けておく}と、本棚に納める順番が前後してもぶつかりません。これがコンフリクト予防の基本的な考え方です。
\end{hintbox}

% ---------------------------------------------------------------------
\section{ブランチの基本操作}

作業ブランチを作って切り替えるところから始めます。

\begin{codebox}[bash]
\begin{code}
# 現在のブランチ一覧を見る（* が今いるブランチ）
git branch

# ブランチを作って、同時にそこへ切り替える
git switch -c feature/add-login

# （参考）古い書き方。同じ意味
git checkout -b feature/add-login
\end{code}
\end{codebox}

\cmd{git switch -c} の \cmd{-c}（create）は「新しく作って切り替える」という意味です。すでにあるブランチへ移動するだけなら \cmd{-c} は不要です。

\begin{codebox}[bash]
\begin{code}
# 既存のブランチへ切り替える
git switch main
git switch feature/add-login

# 不要になったブランチを削除する（マージ済みのものだけ安全に消える）
git branch -d feature/add-login
\end{code}
\end{codebox}

作業ブランチに切り替えたら、あとは\textbf{第1章の黄金フローと同じ}です。「編集 → \cmd{git add} → \cmd{git commit}」を、このブランチの上で繰り返します。

\begin{hintbox}
\textbf{ブランチ名の付け方}：厳密な決まりはありませんが、\cmd{種類/内容} の形が広く使われます。例：新機能なら \cmd{feature/add-login}、バグ修正なら \cmd{fix/typo-in-readme}。目的がひと目で分かる名前にします。
\end{hintbox}

\begin{warnbox}
作業を始める前に、\textbf{\cmd{main} を最新にしてから枝分かれ}するのが安全です（\cmd{git switch main} → \cmd{git pull} → \cmd{git switch -c 新ブランチ}）。古い \cmd{main} から分岐すると、後でコンフリクトが起きやすくなります。
\end{warnbox}

% ---------------------------------------------------------------------
\section{ブランチをリモートへ上げる}

作った作業ブランチは、ローカルにあるだけではチームに共有されません。第2章の \cmd{push} で\textbf{リモートにも同じブランチを作り}ます。

\begin{codebox}[bash]
\begin{code}
# 初回：リモートに同名のブランチを作り、宛先を記憶させる
git push -u origin feature/add-login

# 2回目以降はこれだけでOK
git push
\end{code}
\end{codebox}

\cmd{-u}（\cmd{-{}-set-upstream}）は第2章と同じで、「このローカルブランチの写し先はリモートのこのブランチだ」と結び付けるオプションです。一度付ければ、以降は \cmd{git push} / \cmd{git pull} だけで、今いるブランチが同期されます（図\ref{fig:push-branch}）。

\begin{numfig}{作業ブランチをリモートにも作って同期する}{fig:push-branch}
\begin{tikzpicture}
  \node[flownode, fill=gitOrange, text width=2.8cm] (LB)
        {ローカル\\feature/add-login};
  \node[flownode, fill=gitBlue, text width=2.8cm, right=8cm of LB] (RB)
        {GitHub\\origin/feature/add-login};
  \draw[flowarrow] (LB) -- (RB)
        node[edgelabel, midway, above] {git push -u origin feature/add-login};
\end{tikzpicture}
\end{numfig}

\begin{hintbox}
リモートに上げたブランチは、GitHubのリポジトリ画面のブランチ切り替えメニューから確認できます。これで、他のメンバーもあなたのブランチを見られるようになります。
\end{hintbox}

% ---------------------------------------------------------------------
\section{プルリクエスト（PR）を作る}

\textbf{プルリクエスト（Pull Request、PR）}は、「自分の作業ブランチの変更を \cmd{main} に取り込んでください」という\textbf{依頼}です。ただの合流ではなく、\textbf{レビュー（変更内容の確認）を挟む入口}になります（図\ref{fig:pr-flow}）。

\begin{numfig}{PRは push からマージまでにレビューを挟む}{fig:pr-flow}
\begin{tikzpicture}
  \node[flownode, fill=gitOrange, text width=2.4cm] (B) {feature を push};
  \node[flownode, fill=gitPurple, right=2.2cm of B] (P) {PRを作成};
  \node[flownode, fill=gitBlue, right=2.2cm of P, text width=2.2cm] (R) {レビュー\\（差分を確認）};
  \node[flownode, fill=gitGreen, right=2.2cm of R] (M) {main へマージ};
  \draw[flowarrow] (B) -- (P);
  \draw[flowarrow] (P) -- (R);
  \draw[flowarrow] (R) -- (M);
\end{tikzpicture}
\end{numfig}

\noindent\textbf{PRの作成手順}

\begin{enumerate}[leftmargin=1.6em]
  \item ブランチを \cmd{push} した後、GitHubのリポジトリ画面を開きます。上部に \textbf{「Compare \& pull request」} ボタンが出るのでクリックします（出ない場合は \textbf{Pull requests} タブ → \textbf{New pull request}）。
  \item \textbf{base}（取り込み先＝通常 \cmd{main} または \cmd{develop}）と \textbf{compare}（自分の作業ブランチ）が正しいか確認します。
  \item \textbf{タイトル}と\textbf{説明}を書きます。「何を・なぜ変えたか」を簡潔に書くとレビューが早く進みます。
  \item \textbf{Create pull request} をクリックします。
\end{enumerate}

\screenshot[0.62]{03_branch・PR/PR_Create_Window.png}{PR作成画面（Compare \& pull request）}
\screenshot[0.62]{03_branch・PR/Create_PR_Setting.png}{base と compare を確認し、タイトル・説明を記入して作成する}

\begin{hintbox}
PRを作ると、以降その作業ブランチに \cmd{push} するたびに\textbf{変更が自動でPRに追加}されます。レビューの指摘を直したら、同じブランチで \cmd{add} → \cmd{commit} → \cmd{push} すればOKです。
\end{hintbox}

% ---------------------------------------------------------------------
\section{レビューとマージ}

PRの画面では、\textbf{Files changed} タブで変更の差分（追加＝緑、削除＝赤）を確認できます。チーム開発では、ここで他のメンバーがコメントや承認（Approve）を行います。

PRの差分を見て、\textbf{「この変更で問題ないか」}を確認します。必要に応じてコメントを書き、修正が必要なら作業ブランチに戻って直して \cmd{push} します。レビューが済んだら、PR画面の \textbf{「Merge pull request」} ボタンで \cmd{main} に取り込みます。

\screenshot[0.62]{03_branch・PR/PR_Change_Program.png}{変更内容の確認とチェック}
\screenshot[0.62]{03_branch・PR/Review_Merge.png}{レビューとマージ}

問題がなければ、PR画面の \textbf{「Merge pull request」} ボタンで \cmd{main} に取り込みます。マージが済んだら、次の後片付けをします。

\begin{codebox}[bash]
\begin{code}
# 1. ローカルの main に戻る
git switch main

# 2. マージされた最新の main を取り込む
git pull

# 3. 役目を終えた作業ブランチを削除する（ローカル）
git branch -d feature/add-login

# 4. リモートの作業ブランチも削除する（GitHubのPR画面のボタンでも可）
git push origin --delete feature/add-login
\end{code}
\end{codebox}

\begin{hintbox}
マージ後にブランチを消すのは、\textbf{散らかり防止}のためです。GitHubのPR画面にも、マージ後に \textbf{「Delete branch」} ボタンが出るので、そこから消しても構いません。次のタスクを始めるときは、また \cmd{main} を最新にしてから新しいブランチを切ります。この「\textbf{最新の \cmd{main} から枝を出す}」を守るのが、コンフリクトを減らす一番のコツです。
\end{hintbox}

% ---------------------------------------------------------------------
\section{コンフリクト（競合）を解決する}

\textbf{コンフリクト（競合）}は、\textbf{同じファイルの同じ箇所を、2つの流れが別々に変更した}ときに起こります。Gitが「どちらを採用すべきか自動で判断できない」ため、あなたに解決を求めてきます。

\subsection{なぜ起きるのか}

たとえば、共有リポジトリを元にAさんとBさんが別々に作業し、\textbf{同じ行を違う内容に}書き換えたとします。片方が先にマージした後、もう片方が \cmd{pull} や \cmd{merge} をすると、履歴が食い違ってコンフリクトになります（図\ref{fig:conflict-why}）。

\begin{numfig}{同じ行を別々に変更した2つの流れを合流させるとコンフリクトになる}{fig:conflict-why}
\begin{tikzpicture}
  \node[flownode, fill=gitBlue, text width=3cm] (S) {共有の出発点\\著者: 山田};
  \node[flownode, fill=gitPurple, text width=3.2cm, below left=1.3cm and 0.3cm of S] (A)
        {Aさんが変更\\著者: 山田, 佐藤, 鈴木};
  \node[flownode, fill=gitOrange, text width=3.2cm, below right=1.3cm and 0.3cm of S] (B)
        {Bさんが変更\\著者: 山田, 田中};
  \node[decisionnode, fill=gitAmber, text width=3cm, below=3.4cm of S] (M)
        {同じ行が2通りに…\\どちらが正しい？};
  \node[flownode, fill=gitRed, below=1.2cm of M] (C) {コンフリクト発生};
  \draw[flowarrow] (S) -- (A);
  \draw[flowarrow] (S) -- (B);
  \draw[flowarrow] (A) -- (M);
  \draw[flowarrow] (B) -- (M);
  \draw[flowarrow] (M) -- (C);
\end{tikzpicture}
\end{numfig}

\begin{hintbox}
\textbf{予防が最善です。} 「1タスク＝1ブランチ」「作業前に \cmd{main} を \cmd{pull}」「触るファイルをタスクごとに分ける」を守ると、コンフリクトは大きく減ります。ただしLaTeXの原稿など、同じファイルを複数人で触る場面では起きやすいので、直し方を知っておきましょう。
\end{hintbox}

\subsection{発生したときの見え方}

\cmd{git pull} や \cmd{git merge} の際に、次のようなメッセージが出ます。

\begin{codebox}[text]
\begin{code}
Auto-merging report.tex
CONFLICT (content): Merge conflict in report.tex
Automatic merge failed; fix conflicts and then commit the result.
\end{code}
\end{codebox}

\cmd{git status} を打つと、\cmd{Unmerged paths} として以下のように競合したファイルが表示されます。

\begin{codebox}[text]
\begin{code}
hoge@n171h147 git-practice % git status
On branch main
Your branch is up to date with 'origin/main'.

You have unmerged paths.
  (fix conflicts and run "git commit")
  (use "git merge --abort" to abort the merge)

Unmerged paths:
  (use "git add <file>..." to mark resolution)
        both added:      hoge1.py

no changes added to commit (use "git add" and/or "git commit -a")
\end{code}
\end{codebox}

\subsection{マーカーの読み方}

競合したファイルを開くと、ぶつかった箇所が次のようなマーカーで囲まれています。ここでは \cmd{report.tex} の著者行を、\textbf{自分は「山田・佐藤・鈴木」}、\textbf{相手は「山田・田中」}に変えてぶつかった、という例で見ていきます。

\noindent\textbf{① 発生時の実際の見え方}

\begin{codebox}[text]
\begin{code}
<<<<<<< HEAD
\author{山田, 佐藤, 鈴木}
=======
\author{山田, 田中}
>>>>>>> origin/main
\end{code}
\end{codebox}

\begin{itemize}[leftmargin=1.4em]
  \item \cmd{<<<<<<< HEAD} 〜 \cmd{=======} … \textbf{自分側}（今いるブランチ）の内容 → \cmd{\textbackslash author\{山田, 佐藤, 鈴木\}}
  \item \cmd{=======} 〜 \cmd{>>>>>>> origin/main} … \textbf{相手側}（取り込もうとした内容）→ \cmd{\textbackslash author\{山田, 田中\}}
\end{itemize}

\cmd{HEAD} は「今の自分」を指す目印だと覚えておくと、上が自分・下が相手だと迷いません。

\noindent\textbf{② 手で編集する}

\textbf{どちらを残すか（または両方を活かすか）を決め}、\cmd{<<<<<<<} / \cmd{=======} / \cmd{>>>>>>>} の\textbf{3行をすべて削除}して、正しい1行の状態にします。

\noindent\textbf{③ 解決後の完成形}

たとえば「両方を活かす」なら、最終的にファイルはこうなります。

\begin{codebox}[text]
\begin{code}
\author{山田, 佐藤, 鈴木, 田中}
\end{code}
\end{codebox}

マーカーは1文字も残さず、\textbf{そのまま原稿として通る状態}にするのがゴールです。

\subsection{解決の手順}

全体の流れは次の通りです。\textbf{発生 → 中身を見て「どれを残すか」決める → マーカーを消して正しい状態に → \cmd{add} → \cmd{commit}}。決めるところで、図\ref{fig:conflict-resolve}のように3択に枝分かれします。

\begin{numfig}{コンフリクト解決の流れ（3択に枝分かれ）}{fig:conflict-resolve}
\begin{tikzpicture}[node distance=0.5cm and 1.4cm]
  \node[flownode, fill=gitRed, text width=2cm] (A) {コンフリクト発生};
  \node[decisionnode, fill=gitAmber, text width=2.2cm, right=of A] (C) {どれを残すか決める};
  \node[flownode, fill=gitPurple, text width=2.6cm, right=of C] (C2) {相手側\\山田, 田中};
  \node[flownode, fill=gitBlue, text width=2.6cm, above=of C2] (C1) {自分側\\山田, 佐藤, 鈴木};
  \node[flownode, fill=gitTeal, text width=2.6cm, below=of C2] (C3) {両方\\山田, 佐藤, 鈴木, 田中};
  \node[flownode, fill=gitOrange, text width=2.2cm, right=of C2] (D) {マーカー3行を削除};
  \node[flownode, fill=gitGreen, right=of D] (E) {git add → commit};
  \draw[flowarrow] (A) -- (C);
  \draw[flowarrow] (C) -- (C1);
  \draw[flowarrow] (C) -- (C2);
  \draw[flowarrow] (C) -- (C3);
  \draw[flowarrow] (C1) -- (D);
  \draw[flowarrow] (C2) -- (D);
  \draw[flowarrow] (C3) -- (D);
  \draw[flowarrow] (D) -- (E);
\end{tikzpicture}
\end{numfig}

解決のやり方は2通りあります。チーム開発ではコンフリクトはPR上で表面化するので、\textbf{PR画面のまま直せる方法A（GitHub上）を基本}にしましょう。手元でじっくり直したいときや、ブラウザで扱えない複雑な競合のときは方法B（VS Code）を使います。

\subsubsection{方法A：GitHub（ブラウザ）で解決する（推奨）}

PRを作った後にコンフリクトが起きると、PR画面に \textbf{「This branch has conflicts that must be resolved」} という警告と \textbf{「Resolve conflicts」} ボタンが表示されます。ローカルに戻らず、\textbf{ブラウザだけ}で直せる手軽な方法で、チーム開発ではこれが一番よく使います。

\begin{enumerate}[leftmargin=1.6em]
  \item PR画面の \textbf{「Resolve conflicts」} ボタンを押す。競合したファイルがブラウザ上のエディタで開きます。
\end{enumerate}
\screenshot[0.62]{03_branch・PR/conflicts_first_window.png}{コンフリクト発生時のPR画面}

\begin{enumerate}[leftmargin=1.6em, start=2]
  \item \textbf{\cmd{<<<<<<<} / \cmd{=======} / \cmd{>>>>>>>} のマーカー}が表示されるので、\textbf{正しい状態に直し、マーカー3行をすべて削除}します。
\end{enumerate}
\screenshot[0.62]{03_branch・PR/conflict_fix_file.png}{コンフリクト解決の編集画面}

\begin{enumerate}[leftmargin=1.6em, start=3]
  \item ファイル上部の \textbf{「Mark as resolved」} を押す。競合ファイルが複数あれば、すべて同様に処理します。
\end{enumerate}
\screenshot[0.62]{03_branch・PR/resolve_conflict.png}{コンフリクト解決後のマーク}

\begin{enumerate}[leftmargin=1.6em, start=4]
  \item すべて解決したら \textbf{「Commit merge」} を押して確定します。これで競合が解消され、PRをマージできる状態になります。
\end{enumerate}
\screenshot[0.62]{03_branch・PR/commit_merge.png}{コンフリクト解決後のコミット}

\begin{warnbox}
\textbf{「Resolve conflicts」ボタンが出ないことがあります。} 競合が複雑すぎる場合、GitHubはブラウザでの解決を許可しません。そのときは方法B（VS Code）で、ローカルに取り込んでから解決します。
\end{warnbox}

\subsubsection{方法B：VS Codeで解決する}

VS Codeは競合を検知すると、専用のボタンで解決を手伝ってくれます。\textbf{マーカーを自分で消す必要はありません。}

\begin{enumerate}[leftmargin=1.6em]
  \item \textbf{競合したファイルを開く。} サイドバーの \textbf{ソース管理（Source Control）} パネルに、競合ファイルが \textbf{「Merge Changes」} としてまとめて表示されます。エクスプローラー上では、そのファイル名の後ろに \cmd{!}（Conflict）マークが付きます。
\end{enumerate}
\screenshot[0.62]{03_branch・PR/vscode_conflict_file.png}{VS Codeでのコンフリクト画面}

\begin{enumerate}[leftmargin=1.6em, start=2]
  \item \textbf{ファイルを開くと、「マージエディターで解決」ボタンが表示されます。} この状態で \cmd{<<<<<<<} / \cmd{=======} / \cmd{>>>>>>>} のマーカーを削除し、正しい内容に書き換えることもできます。
\end{enumerate}
\screenshot[0.62]{03_branch・PR/vscode_open_conflict_file.png}{VS Codeでのマージエディター起動}

\begin{enumerate}[leftmargin=1.6em, start=3]
  \item VS Codeのマージエディターを使うと、\textbf{左＝自分側／右＝相手側／下＝結果}の3ペインで、どちらを採用するかをチェックボックスで選べます。両方を活かして手直ししたいときは、下段を直接編集します。完了したら左側のソース管理パネルで \cmd{+}（ステージ）を押して、コミットすれば解決完了です（コマンドの \cmd{git add} → \cmd{git commit} と同じ操作です）。
\end{enumerate}
\screenshot[0.62]{03_branch・PR/vscode_merge_editor.png}{VS Codeのマージエディター（3ペイン）}

どちらの方法でも、複数ファイルで起きた場合は\textbf{すべてのファイルを解決してから}マージを確定します。

\begin{hintbox}
\textbf{コマンドラインだけで解決したい場合}は、ファイルを直接開いてマーカーを消し、\cmd{git add <ファイル>} → \cmd{git commit} で確定します。方法A・Bのボタン操作も、内部ではこれと同じことをしています。
\begin{codebox}[bash]
\begin{code}
# 解決したファイルを「解決済み」として登録する
git add report.tex

# コミットして競合の解決を確定する（メッセージは自動で用意される）
git commit
\end{code}
\end{codebox}
\end{hintbox}

\begin{warnbox}
途中で「やっぱりやめたい」ときは \cmd{git merge -{}-abort} で、マージ開始前の状態に戻せます。慌てて壊す前に、まず落ち着いてこのコマンドを思い出してください。
\end{warnbox}

% ---------------------------------------------------------------------
\section{実践フロー：ブランチ運用の一周}

ここまでを1つの流れにまとめます。\textbf{この一周を繰り返す}のが、チーム開発の基本リズムです（図\ref{fig:branch-cycle}）。

\begin{numfig}{ブランチ運用の一周}{fig:branch-cycle}
\begin{tikzpicture}[node distance=0.75cm]
  \node[flownode, fill=gitGreen, text width=7cm] (A) {git switch main → git pull（最新化）};
  \node[flownode, fill=gitOrange, text width=7cm, below=of A] (B) {git switch -c feature/xxx（枝を出す）};
  \node[flownode, fill=gitRed, text width=7cm, below=of B] (C) {編集 → git add → git commit};
  \node[flownode, fill=gitBlue, text width=7cm, below=of C] (D) {git push -u origin feature/xxx};
  \node[flownode, fill=gitPurple, text width=7cm, below=of D] (E) {GitHubでPR作成};
  \node[flownode, fill=gitBlue, text width=7cm, below=of E] (F) {レビュー → マージ};
  \node[flownode, fill=gitGreen, text width=7cm, below=of F] (G) {main を pull、済んだブランチを削除};
  \draw[flowarrow] (A) -- (B);
  \draw[flowarrow] (B) -- (C);
  \draw[flowarrow] (C) -- (D);
  \draw[flowarrow] (D) -- (E);
  \draw[flowarrow] (E) -- (F);
  \draw[flowarrow] (F) -- (G);
  \draw[flowarrow] (G.east) to[out=0, in=0]
        node[edgelabel, midway, right] {繰り返す} (A.east);
\end{tikzpicture}
\end{numfig}

\begin{codebox}[bash]
\begin{code}
# 1. main を最新にする
git switch main
git pull

# 2. 作業ブランチを切る
git switch -c feature/add-login

# 3. 編集して記録する（第1章の黄金フロー）
git add .
git commit -m "ログイン機能を追加"

# 4. リモートへ上げる
git push -u origin feature/add-login

# 5. GitHubでPRを作成 → レビュー → マージ（ブラウザで操作）

# 6. 後片付け
git switch main
git pull
git branch -d feature/add-login
\end{code}
\end{codebox}

% ---------------------------------------------------------------------
\section*{この章のまとめ}
\addcontentsline{toc}{section}{この章のまとめ}

\begin{itemize}[leftmargin=1.4em]
  \item \textbf{\cmd{main} は直接触らず}、タスクごとに\textbf{作業ブランチ}を切って編集する
  \item 作業ブランチは \cmd{git push -u origin ブランチ名} でリモートにも共有する
  \item \textbf{PR}は「\cmd{main} に取り込んでほしい」という依頼であり、\textbf{レビューの入口}
  \item マージ後は \textbf{\cmd{main} を \cmd{pull} → 済んだブランチを削除}して片付ける
  \item \textbf{コンフリクト}は同じ箇所の二重変更で起きる。マーカーを消して \cmd{add} → \cmd{commit} で解決
  \item 何より、\textbf{最新の \cmd{main} から枝を出す}ことがコンフリクト予防の要
\end{itemize}

これで、ローカル（第1章）→ GitHub連携（第2章）→ チーム開発（第3章）の一通りが揃いました。日々の作業でコマンドを忘れたら第5章「逆引きチートシート」を、エラーやトラブルに出くわしたら第4章「トラブルシューティング編」を、それぞれ辞書代わりに使ってください。
